\section{Validation information and feasible inference}\label{sec:inference}
\subsection{An existing probability sample of items}\label{sec:audit}
Suppose predictions are frozen before items are selected for verification.
The audit observes $y_j$ for selected items, with known inclusion
probabilities. This is one source of inference about the complete-label
panel coefficient. It does not require optimal selection or calibrated
scores, but it does require correctly verified labels and a known design.
Audit-based regression correction predates current surrogate-assisted
methods \citep{shepherd2012audit}; survey calibration and design-based
supervised learning provide related estimating-function constructions
\citep{lumley2011calibration,egami2023}.

For a linear target, let $R_j$ indicate verification and let
$\pi_j=\Pr(R_j=1)>0$. The difference estimator is
\begin{equation}\label{eq:mainauditlinear}
\widetilde\theta=w^{\mathsf T}q+
             \sum_j\frac{R_j}{\pi_j}w_j(y_j-q_j).
\end{equation}
It is design unbiased for $w^{\mathsf T}y$. For several targets, replace
$w_j$ by its vector of weights across targets. This preserves the joint
correction contributed by one verified item. The exact design covariance
uses the first- and second-order inclusion probabilities, not a separate
error draw for each downstream record (Appendix~\ref{app:validation}).

\subsection{Regression moments and their joint covariance}
For a constructed covariate, equation~\eqref{eq:truecovariate} has a linear
numerator and quadratic denominator. If both regression sides are
constructed, both moments are quadratic. More generally, each component of
the normal equations has form
$m_s(y)=a_s^{\mathsf T}y+y^{\mathsf T}G_sy$, with $G_s$ symmetric.
Writing $r=y-q$ gives
\begin{equation}\label{eq:polynomialaudit}
\boldsymbol m(y)=\boldsymbol m(q)+\sum_{A\in\mathcal A}c_A,
\qquad
\widehat{\boldsymbol m}=\boldsymbol m(q)+
                 \sum_{A\in\mathcal A}\frac{R_A}{p_A}c_A,
\end{equation}
where $\mathcal A$ contains singleton items and unordered item pairs,
$R_A=\prod_{j\in A}R_j$, and $p_A=\Pr(R_A=1)>0$. The $s$th singleton
contribution is $r_j\{a_{sj}+2(G_sq)_j\}+G_{s,jj}r_j^2$; the pair
contribution for $j<\ell$ is $2G_{s,j\ell}r_jr_\ell$. Both are observed
whenever their items are verified. Quadratic finite-population totals can
be treated as totals on such a derived population \citep{sitter2002quadratic}.

\begin{proposition}[Joint audit moments]\label{prop:momentcovariance}
Conditional on the panel and frozen predictions,
$\E_R[\widehat{\boldsymbol m}]=\boldsymbol m(y)$. If every relevant
$p_{A\cup B}$ is known and positive, an unbiased estimator of its covariance is
\begin{equation}\label{eq:polynomialvariance}
\widehat V=
\sum_{A,B\in\mathcal A}
\frac{R_{A\cup B}(p_{A\cup B}-p_Ap_B)}
     {p_Ap_Bp_{A\cup B}}\,c_Ac_B^{\mathsf T}.
\end{equation}
Thus quadratic regression moments require inclusion probabilities for
sets of up to four items.
\end{proposition}
\begin{proof}
Each term in equation~\eqref{eq:polynomialaudit} has expectation $c_A$.
The covariance of the $A$ and $B$ terms is
$(p_{A\cup B}-p_Ap_B)c_Ac_B^{\mathsf T}/(p_Ap_B)$.
Multiplication by $R_{A\cup B}/p_{A\cup B}$ gives an unbiased estimate of
each term.
\end{proof}
This is a Horvitz--Thompson covariance calculation, applied jointly to the
regression moments. It does not introduce a new regression estimator.
Independent Bernoulli item sampling gives a particularly simple
implementation because disjoint item sets have zero covariance;
Appendix~\ref{app:momentvariance} gives the expression used below.

For $\beta=N/D$, the gradient of the regression map is
$g=(1/D,-N/D^2)^{\mathsf T}$. A feasible delta SE uses
\begin{equation}\label{eq:auditratiose}
\widehat\beta=\widehat N/\widehat D,\qquad
\widehat{\operatorname{se}}(\widehat\beta)
=\{\widehat g^{\mathsf T}\widehat V\widehat g\}^{1/2}.
\end{equation}
A normal interval is justified when the joint moments have a normal
approximation, the denominator stays sufficiently far from zero, the
linearization remainder is negligible, and the estimated variance is
consistent on the coefficient's variance scale. With several constructed
covariates, the moment vector includes all regressor--outcome and
regressor--regressor products; the derivative of the full regression map
must be applied to their joint covariance.

These conditions do not follow from unbiasedness of either moments or
covariance. In a finite audit, $\widehat D$ can be nonpositive and an
unbiased covariance estimate can be indefinite. Such draws do not support
the stated Wald SE. A few influential items can also leave a nonnormal
coefficient distribution even when the covariance formula is exact.
Section~\ref{sec:simulation} evaluates both SE estimation and coverage,
including these failures.

A finite-sample alternative starts with a confidence set $\mathcal Y_A$
for the complete label vector, such as the simultaneous count certificate
in Appendix~\ref{app:validation}. On its coverage event,
\[
\beta(y)\in\{\beta(z):z\in\mathcal Y_A,\ D(z)>0\}.
\]
The coverage implication holds for any fixed downstream map; it is the
usual projection of a confidence set, also used in prediction-powered
inference \citep{angelopoulos2023ppi}. Computing useful nonlinear endpoints
and handling possible rank failures are additional tasks. The resulting
set may be very wide or unbounded.

\subsection{Calibrated predictions and estimated stages}\label{sec:calibration}
Calibration is not required for fixed-label bounds or the probability
audit above. A score-based population interpretation needs different
assumptions. If information $\mathcal H$ includes the loading matrices and
controls and $p_j=\E[y_j\mid\mathcal H]$, then
$Cp=\E[Cy\mid\mathcal H]$. Under
$Y=V\gamma+\beta Cy+\varepsilon$, with
$\E[\varepsilon\mid\mathcal H]=0$ and a full-rank conditional-mean design,
regressing $Y$ on $(V,Cp)$ identifies the population $\beta$.
Ordinary calibration across items, $\Pr(y=1\mid p=s)=s$, does not establish
these conditional means or the exclusion restriction. Appendix~\ref{app:calibration}
gives a weighted counterexample and separates expected aggregates from
realized finite-item targets.

Even correct conditional means do not supply the covariance of the
remaining errors across records. Estimating calibration adds another
estimation stage. A stacked estimating-equation covariance or a justified
bootstrap must propagate that stage together with the downstream fit
\citep{murphy1985twostep,boe2023calibration,boe2024sandwich}.
Methods that validate complete regression rows and permit nonuniform
sampling offer useful precedents \citep{kluger2025imputed}; the sampled
unit must still match the verified information available here.

If uncertainty instead concerns repeated classifier fitting, the fitting
sample and all item predictions must be varied jointly under a resampling
scheme justified for that process. Independently perturbing each downstream
record invents independent versions of a shared label. No known sampling
probabilities, error model, or training experiment can be recovered from a
convenience confusion table alone.

\subsection{What counts omit}
Exact target-item false-positive and false-negative counts constrain which
labels can be wrong, but discard their locations in the loading matrices.
For linear targets, Proposition~\ref{prop:bounds} gives the sharp interval
obtained by sorting their item weights. Counts suffice within a predicted
class only when all relevant weights are equal, or its error count is zero
or exhaustive (Proposition~\ref{prop:sufficient}). These bounds concern
identification from a summary, not a standard error. Their proof and the
audit count certificate are in Appendices~\ref{sec:identification} and
\ref{app:validation}.
